\subsection{Visit changes, coverage and prize restoration}
The fixed-support replay provides conditional evidence that reordering alone
does not account for all of the observed training advantage. Both alternatives
started from the same saved MDM states and had the corresponding remaining
allowance. The comparison nevertheless excludes generation of those states,
uses previously examined training geometries and evaluates the complete
visit-changing phase. It cannot separately identify the effect of the
destruction size, greedy scoring, temporary infeasibility, acceptance rule or
sequence of LKH calls.

A descriptive check of the 27 completed training trajectories found that
the final visited set differed from the pre-repair set in 23 runs. Removing
the locations absent from the final set, without yet adding the final new
locations, would violate the prize quota in 20 runs and coverage in five.
The five coverage cases were included among the 20 quota cases. These are
feasibility witnesses for net differences between two feasible endpoints.
They do not reconstruct an intermediate destruction state, establish the
necessity of a particular insertion, or prove which search component caused
the final improvement.

The input structure further limits a coverage-based explanation. In gr229,
kroA200, pr226 and rd400, the mandatory locations already cover every demand.
Because all admissible tours retain those locations, each coverage constraint
is automatically satisfied in these four instances. Optional-visit changes
there still affect routing and prize collection, but their gains cannot be
attributed to restoring missing demand coverage. In all nine selected training
instances, mandatory prize alone is insufficient. Thus the available evidence
supports examining prize and visit selection alongside coverage, rather than
assuming that coverage restoration is the dominant mechanism throughout the
training set.

\subsection{Restricted integer repair and computational allocation}
The evaluated joint-insertion model did not improve on greedy repair in the
primary mean comparison. Its advantage over the equal-size random-gap control
suggests that the chosen insertion positions matter within that experiment;
it does not establish an advantage for MIP repair over the simpler repair
routine. Restricting candidates and positions can exclude useful tours, and
a model that restores coverage and prize exactly need not find a lower-cost
tour within its allowance. The present comparison does not isolate which
restriction or runtime component explains the unfavorable mean effect.

The same distinction applies to path contraction and local optimality.
An exact representation of a restricted repair preserves the meaning of its
feasible solutions and objective values. It does not show that solving the
restricted model is a good allocation of an end-to-end search budget. Likewise,
an optimal restricted repair is not a certificate for the unrestricted PCCTP.
The negative repair result is retained separately from the uniform-removal,
greedy-repair hybrid; the successful training screen of the latter cannot be
used as evidence for an integer component that it did not execute.

\subsection{Limits of the current performance evidence}
The training set contains nine selected difficult geometries, rather than a
representative sample of the full published evaluation set. The seven
validation geometries have also been reused. Their two positive means and
zero median, together with the failed breadth criterion, limit any claim of
general improvement. The twelve reserved test geometries remain unopened;
their existence is not evidence of performance on them. Descriptive omission
analyses do not repair selection bias or supply independent confirmation.

Search horizons create another limitation. A short-budget comparison can end
before MDM activates its mining component. It then compares the observed
algorithm trajectories within that horizon, rather than isolating the value
of mining or outperforming a completed 70-iteration reference. The completed
300-second experiment favored the 60-second initialization variant in the
selected training set, while the 10-second variant failed its screen.
This difference cautions against treating additional repair time as a
substitute for adequate initialization. It does not isolate the causal
effect of initialization length, because the resulting search trajectories
and computational allocation also differ. The planned iteration-limited
comparisons address a different stopping rule and remain necessary.
The reference-defined duration
rule also requires disclosure of the reference-first order and possible
machine-state drift within a pair.

Finally, the recovered baseline includes declared decoder and elite-count
corrections, and measured times refer to the present hardware and executable
versions. They do not establish speedups over seconds reported on the older
processor. Any final performance contribution requires valid complete
comparisons, a separately frozen broader evaluation and an account of failures,
mining exposure and computational cost. The completed results currently justify
continued evaluation of the visit-changing hybrid, while retaining the failed
validation screen and negative integer-repair comparison as constraints on
its interpretation.
